\chapter{Osteoporosis}
\index{Osteoporosis}
\label{sec:Osteoporosis}

\section{Overview}
It affects approximately 200 million people worldwide, predominantly postmenopausal women and the elderly. This metabolic disorder involves dysregulation of normal bodily processes, often requiring ongoing monitoring and management. It is increasingly common due to lifestyle and dietary factors. This degenerative condition involves progressive deterioration of affected tissues, leading to gradual loss of function. It is more common with advancing age. Metabolic disorders involve disruption of normal biochemical processes essential for health. These conditions may result from enzyme deficiencies, hormonal imbalances, or lifestyle factors. Many metabolic disorders are chronic and require ongoing management. Lifestyle modifications are often the cornerstone of treatment.
\section{Causes}
The underlying mechanisms involve complex interactions between genetic predisposition and environmental factors. Disruption of normal physiological processes leads to the characteristic features of the condition. Multiple contributing factors may act synergistically to trigger or worsen the condition.
\section{Risk Factors}


\begin{itemize}[noitemsep]
  \item Family history
  \item Obesity and sedentary lifestyle
  \item Poor dietary habits
  \item Advancing age
  \item Hormonal imbalances
  \item Certain medications
  \item Genetic predisposition
  \item Repetitive use or overuse (joints)
  \item Previous injuries
  \item Obesity
  \item Smoking and alcohol use
  \item Genetic predisposition and family history
  \item Lifestyle factors (diet, exercise, smoking, alcohol)
  \item Environmental exposures
  \item Underlying medical conditions
  \item Medication use
\end{itemize}
\section{Signs and Symptoms}

\subsection{Common symptoms}
\begin{itemize}[noitemsep]
  \item You may experience fragility fractures commonly at the hip, vertebrae (compression fractures), and distal radius (Colles' fracture).
  \item Pain or discomfort in the affected area
  \end{itemize}

\subsection{Emergency warning signs}
\begin{itemize}[noitemsep]
  \item Severe or worsening symptoms of osteoporosis require immediate medical evaluation
\end{itemize}
\section{Progression and Stages}
Osteoporosis is a systemic skeletal disease characterized by low bone mineral density (BMD) and deterioration of bone microarchitecture, resulting in increased bone fragility and fracture risk. The condition may progress through different stages of severity over time. Early stages often involve mild or subtle symptoms that may be overlooked. Moderate stages involve more noticeable symptoms affecting daily function. Advanced stages may involve significant impairment and complications.
\section{Complications}
\begin{itemize}[noitemsep]
  \item Disease progression leading to worsening symptoms
  \item Secondary infections or injuries
  \item Side effects of treatment
  \item Reduced quality of life
  \item Emotional and psychological impact
\end{itemize}
\section{Diagnosis}
\begin{itemize}[noitemsep]
  \item Doctors diagnose this based on dual-energy X-ray absorptiometry (DXA): T-score $\leq$-2.5 at the hip or spine defines osteoporosis
  \item T-score between -1 and -2.5 is osteopenia. The FRAX tool estimates 10-year fracture risk.
  \item Comprehensive medical history and physical examination
  \item Laboratory tests to assess organ function
  \item Imaging studies to visualize affected structures
  \item Specialized testing depending on the affected system
  \item Consultation with relevant specialists
\end{itemize}
\section{Prevention}
\begin{itemize}[noitemsep]
  \item Outlook is improved with early detection and treatment
  \item Fracture risk is significantly reduced with appropriate therapy.
  \item Maintain a healthy lifestyle with balanced diet and exercise
  \item Avoid known risk factors
  \item Attend regular medical check-ups for early detection
  \item Follow recommended screening guidelines
\end{itemize}
\section{Treatment Options}
Treatment options include calcium (1000--1200 mg/day) and vitamin D (800--1000 IU/day), weight-bearing exercise, bisphosphonates (alendronate, zoledronic acid, first-line), denosumab (RANKL inhibitor), teriparatide and abaloparatide (parathyroid hormone analogs), and romosozumab (anti-sclerostin antibody). Dietary modifications and nutritional counseling Pharmacotherapy to correct metabolic abnormalities Hormone replacement therapy when indicated Regular monitoring of metabolic parameters Weight management programs Exercise prescription Management of associated conditions Patient education for self-management

\begin{itemize}[noitemsep]
  \item Medications to manage symptoms and address underlying mechanisms
  \item Lifestyle modifications and self-care measures
  \item Physical therapy and rehabilitation
  \item Surgical intervention when indicated
  \item Regular monitoring and follow-up care
\end{itemize}
\section{Living with It}
\begin{itemize}[noitemsep]
  \item Follow your treatment plan as prescribed
  \item Maintain regular follow-up with your healthcare provider
  \item Make necessary lifestyle adjustments
  \item Seek support from family, friends, or support groups
  \item Stay informed about your condition
\end{itemize}
